\begin{textsample}{NEG Dim 4 – 2022\_c5 – Score -1.00 – 2022\_c5/enade\_2022\_publi\_d731a0a567.txt}  \label{ex:f4_neg_074}
CAPÍTULO 6 DISTRIBUIÇÃO DE CONCEITOS A finalidade deste capítulo é apresentar um panorama nacional sobre a distribuição do Conceito Enade 2022.
Para traçar esse panorama recorre -se a tabelas e gráficos, que são acompanhados de comentários e análises que contribuem para a contextualização dos dados neles expostos.
Para visualizar o panorama do Enade/2022, em nível nacional, na primeira seção do capítulo, os dados sobre o número de cursos participantes são organizados por grande região.
Na segunda, o mapeamento do conceito do exame é feito por categoria administrativa das IES e modalidade de oferta dos cursos, de maneira cruzada às grandes regiões.
Na última seção, a lógica de apresentação dos dados é feita por organização acadêmica e, também, grande região de funcionamento dos cursos, no país como um todo. 6.1 PANORAMA NACIONAL DA DISTRIBUIÇÃO DOS CONCEITOS A lógica adotada na apresentação dos dados, nesta seção e nas duas seguintes, pressupõe o registro de duas informações de natureza reiterativa.
O primeiro diz respeito aos valores possíveis do Enade, que variam em níveis de 1 a 5, sendo que à medida que esses valores aumentam, o curso é melhor avaliado e ao fato de que aqueles cursos com somente um ou sem nenhum concluinte participante não obtêm o Conceito Enade, ficando Sem Conceito ( SC ), como ressaltado na Apresentação desse relatório.
O segundo registro é que a diferença entre os cursos tabulados neste capítulo e no Capítulo 3 corresponde aos cursos que não foram avaliados, ou seja, àqueles, em princípio, sem estudantes concluintes que tivessem sido inscritos para a prova.
Na Tabela 6.1, são apresentadas a quantidade e a distribuição de cursos de Publicidade e Propaganda participantes do Enade/2022, por faixa de conceito e grande região.
Observando -se os dados da Tabela 6.1, nota -se que, dos 345 cursos participantes, 154 ( 44,6 \% ) classificaram -se com conceito 3, o valor modal.
Esse foi também o conceito modal em todas as regiões : Nordeste ( 39,1 \% ), Norte ( 57,1 \% ), Sudeste ( 45,7 \% ), Sul ( 40,6 \% ) e Centro-Oeste ( 50,0 \% ).
O conceito 4 foi o segundo mais frequente em nível nacional ( 24,6 \%, correspondendo a 85 cursos ), tendo sido o segundo conceito mais frequente na maioria das regiões, exceto no Nordeste.
O conceito 2 foi o terceiro ( 13,6 \%, correspondendo a 47 cursos ).
Houve, ainda, 29 cursos ( 8,4 \% ) que receberam o conceito 5, e 7 cursos ( 2,0 \% ) que receberam o conceito 1.
Dos 345 cursos de Publicidade e Propaganda, 23 ( 6,7 \% ) ficaram Sem Conceito ( SC ).
A região Norte participou com 14 cursos, o correspondente a 4,1 \% do total nacional.
Desses, oito cursos receberam o conceito 3, o que equivale a 57,1 \% do total regional, o conceito modal na região, como já comentado.
A um curso ( 7,1 \% ) atribuiu -se o conceito 2, a dois cursos ( 14,3 \% ), o conceito 4 e, a um curso ( 7,1 \% ), o conceito 5.
Dois cursos ficaram Sem Conceito ( SC ) e nenhum obteve conceito 1.
A região Nordeste participou com 46 cursos, o correspondente a 13,3 \% do total nacional.
Nessa região, 18 cursos ( 39,1 \% em termos regionais ) obtiveram conceito 3, o conceito modal para a região, como já comentado.
Os conceitos 1 e 2 foram atribuídos, respectivamente, a 1 e a 10 cursos ( 2,2 \% e 21,7 \%, respectivamente ).
Ainda nessa região, 9 cursos ( 19,6 \% ) foram avaliados com o conceito 4, e quatro cursos ( 8,7 \% ) foram avaliados com o conceito 5.
Outros quatro cursos ficaram Sem Conceito ( SC ).
Para os 188 cursos participantes da região Sudeste ( 54,5 \% do total nacional ), o conceito 3 foi o conceito modal, como já anteriormente assinalado, tendo sido obtido por 86 cursos ( 45,7 \% ).
O conceito 1 foi atribuído a seis cursos ( 3,2 \% ), e o conceito 2, a 28 cursos ( 14,9 \% ).
Receberam os conceitos 4 e 5, respectivamente, 44 cursos ( 23,4 \% ) e 13 cursos ( 6,9 \% ).
Nessa região, 11 cursos ficaram Sem Conceito ( SC ).
Os 69 cursos da região Sul corresponderam a 20,0 \% do total nacional.
Como já indicado, houve predominância do conceito 3 ( conceito modal ), atribuído a 40,6 \% dos cursos da região, ou seja, atribuído a 28 dos 69 cursos participantes na região Sul.
O conceito 2 foi atribuído a seis cursos ( 8,7 \% ) ; o conceito 4, a 21 cursos ( 30,4 \% ) ; e o conceito 5, a nove cursos ( 13,0 \% ).
Na região Sul, cinco cursos ficaram Sem Conceito ( SC ) e nenhum obteve conceito 1.
Já dos 28 cursos participantes na região Centro-Oeste ( 8,1 \% do total nacional ), 14 ( 50,0 \% em termos regionais ) receberam conceito 3, o conceito modal, como já destacado.
Dois cursos ( 7,1 \% ) obtiveram o conceito 2, nove cursos ( 32,1 \% ) obtiveram o conceito 4, e outros dois ( 7,1 \% ), o conceito 5.
Nessa região, nenhum curso obteve o conceito 1 e apenas um ficou Sem Conceito ( SC ).
Em um gráfico de distribuição cumulativa, em cada valor de abscissa ( eixo x ), o valor apresentado no eixo das ordenadas ( eixo y ) é o valor acumulado da variável ( Conceito Enade ).
Em linhas gerais, em um gráfico desta mesma tipologia, poligonais mais à esquerda representam, no caso dos dados apresentados nesta seção, as grandes regiões com a pior distribuição, e poligonais mais à direita, as grandes regiões com os melhores conceitos.
A partir destas diretrizes, no Gráfico 6.1, é apresentada a distribuição cumulativa do Conceito Enade, segundo a grande região.
Nesse gráfico, poligonais mais à esquerda representam Conceito Enade mais baixos, e poligonais mais à direita, Conceito Enade mais altos. 6.2 CONCEITOS POR CATEGORIA ADMINISTRATIVA, MODALIDADE DE OFERTA DOS CURSOS E GRANDE REGIÃO Na Tabela 6.2, é apresentada a distribuição dos cursos de Publicidade e Propaganda participantes do Enade/2022, por categoria administrativa e modalidade de oferta, de acordo com os conceitos por eles alcançados, segundo as grandes regiões brasileiras.
De acordo com as informações da Tabela 6.2, dos 345 cursos participantes, 315 ( 91,3 \% ) eram ministrados em instituições privadas, e 30 ( 8,7 \% ), em instituições públicas.
Quanto à modalidade de oferta, 16 cursos eram a distância, e 329, Presenciais.
Ainda conforme as informações da mesma tabela, em termos nacionais, entre instituições públicas e privadas, observa -se que, dos 29 cursos avaliados com conceito 5, 9 foram oferecidos em IES públicas, e 20, em IES privadas.
Dos 315 cursos participantes de IES privadas, o conceito 3 foi o valor modal, atribuído a 147 cursos ( 46,7 \% ).
Dos demais cursos de instituições privadas participantes, sete cursos ( 2,2 \% ) foram avaliados com o conceito 1, 45 cursos ( 14,3 \% ) foram avaliados com o conceito 2, 73 cursos ( 23,2 \% ) foram avaliados com o conceito 4 e, 20 cursos ( 6,4 \% ), com o conceito 5.
Nessa categoria, 23 cursos ( 7,3 \% ) ficaram Sem Conceito ( SC ).
Na rede pública, o conceito modal foi o 4, atribuído a 12 cursos, o correspondente a 40,0 \% dos 30 cursos da categoria.
Dos demais cursos participantes, dois ( 6,7 \% ) receberam o conceito 2, e sete ( 23,3 \% ), o conceito 3.
O conceito 5 foi atribuído a 9 cursos ( 30,0 \% ), como já comentado.
Nessa categoria administrativa, nenhum curso ficou Sem Conceito ( SC ) ou com conceito 1.
Dos 329 cursos participantes da modalidade presencial, o conceito 3 foi o valor modal, atribuído a 144 cursos ( 43,8 \% ).
Dos demais cursos nessa modalidade de oferta, sete cursos ( 2,1 \% ) receberam o conceito 1 ; 46 cursos ( 14,0 \% ) receberam o conceito 2 ; 83 cursos ( 25,2 \% ), o conceito 4 ; e 29 cursos ( 8,8 \% ), o conceito 5.
Nessa modalidade de oferta, 20 cursos ( 6,1 \% ) ficaram Sem Conceito ( SC ).
Dos 16 cursos participantes da modalidade a distância, o conceito 3 foi o valor modal, atribuído a 10 cursos ( 62,5 \% ).
Dos demais cursos nessa modalidade de oferta, 1 curso ( 6,2 \% ) recebeu o conceito 2 ; 2 cursos ( 12,5 \% ) receberam o conceito 4 ; e nenhum curso recebeu conceito 1 ou 5.
Nessa modalidade de oferta, três cursos ( 18,7 \% ) ficaram Sem Conceito ( SC ).
Na análise por região, observa -se, que, na região Norte, as instituições públicas participaram com apenas um curso, o qual obteve conceito 4.
As instituições privadas participaram com 13 cursos na região Norte ( 92,8 \% do total regional ), e dois cursos ficaram Sem Conceito ( SC ).
O conceito modal nas instituições privadas na região foi o 3, atribuído a 8 cursos, o correspondente a 61,5 \% dessa categoria na região.
Os conceitos 2, 4 e 5 foram atribuídos a um curso, cada e, o conceito 1, não foi atribuído a nenhum curso.
Dos cursos oferecidos na região Norte, 10 eram presenciais, e dois ficaram Sem conceito ( SC ).
Dos 10 cursos oferecidos, cinco obtiveram o conceito 3 ( modal ) ; 2, o conceito 4 e um, o conceito 5.
Os conceitos 1 e 2 não foram atribuídos a nenhum curso na região Norte e na modalidade presencial.
Quatro eram na modalidade a distância, tendo três recebido o conceito modal 3 e o outro, o conceito 2.
Nenhum curso ficou Sem Conceito ( SC ) ou recebeu os demais conceitos.
Na região Nordeste, a rede privada concentra 38 dos 46 cursos participantes, o equivalente a 82,6 \% do total da região.
Entre os cursos oferecidos por IES privadas no Nordeste, 17 cursos foram avaliados com conceito 3, o conceito modal.
Um curso obteve o conceito 1 ; nove, o conceito 2 ; cinco, o conceito 4 ; e dois cursos obtiveram o conceito 5.
Nessa categoria administrativa e grande região, quatro cursos ficaram Sem Conceito ( SC ).
As instituições públicas da região Nordeste participaram com 8 cursos ( 17,4 \% do total da região ).
Desses, 4 ( 50,0 \% ) obtiveram conceito 4, o valor modal, seguidos por 2 cursos, que receberam o conceito 5.
Os conceitos 2 e 3 foram atribuídos a 1 curso cada.
Nenhum curso ficou Sem Conceito ( SC ) ou com conceito 1 nessa região.
Dos cursos oferecidos na região Nordeste, todos ( 46 ) eram presenciais : quatro ficaram Sem conceito ( SC ), um obteve com o conceito 1 ; 10, com o conceito 2 ; 18, com o conceito 3, o valor modal ; nove, com o conceito 4 ; e quatro, com o conceito 5.
Na região Sudeste, a proporção de cursos da rede privada, 94,7 \%, correspondeu a 178 dos 188 cursos participantes.
Nessa combinação de categoria administrativa e grande região, o conceito modal foi o 3, atribuído a 83 cursos. 11 cursos ficaram Sem Conceito ( SC ).
Os demais foram avaliados com os seguintes conceitos : 1 ( seis cursos ), 2 ( 27 cursos ), 4 ( 41 cursos ) e 5 ( 10 cursos ).
Dos 10 cursos oferecidos por instituições públicas na região Sudeste, três cursos foram atribuídos aos conceitos 3, 4 e 5, cada.
Um curso obteve conceito 2 e nenhum dos cursos ficou Sem Conceito ( SC ) ou com conceito 1.
Também na região Sudeste, grande parcela dos cursos ( 98,4 \% ) foi oferecida na modalidade presencial : 185 de 188 cursos.
Para esses, a moda foi o conceito 3, com 85 cursos.
Os demais cursos presenciais da região foram \textbf{alocados} nos conceitos 1 ( seis cursos ), 2 ( 28 cursos ), 4 ( 43 cursos ) e 5 ( 13 cursos ).
Ficaram Sem Conceito ( SC ) 10 cursos presenciais na região Sudeste.
Na região Sudeste, apenas 3 cursos foram oferecidos na modalidade a distância, tendo sido um curso com conceito 3 ; outro com conceito 4 ; e um ficou Sem Conceito ( SC ).
As instituições privadas concentraram 61 dos 69 cursos participantes da região Sul, 88,4 \% do total regional.
Desses, 26 cursos obtiveram conceito 3, o conceito modal.
Nessa combinação de categoria administrativa e grande região, seis cursos receberam o conceito 2 ; 18 receberam o conceito 4 ; outros seis, o conceito 5.
Cinco cursos ficaram Sem Conceito ( SC ) e nenhum obteve conceito 1.
As instituições públicas na região Sul participaram com 8 cursos ( 11,6 \% dos cursos da região ), aos quais foram atribuídos os conceitos 3 ( dois cursos ), 4 e 5 ( três cursos cada, modais ) e nenhum curso da região Sul oferecido por IES públicas ficou Sem Conceito ( SC ) ou recebeu os conceitos 1 e 2.
Na região Sul, 61 dos cursos foram oferecidos na modalidade presencial.
Para esses, a moda foi o conceito 3, com 23 cursos.
Os demais cursos presenciais da região foram \textbf{alocados} nos conceitos 2 ( seis cursos ), 4 ( 20 cursos ) e 5 ( nove cursos ).
Ficaram Sem Conceito ( SC ) cinco cursos e nenhum obteve conceito 1.
Na modalidade a distância, foram oferecidos 8 cursos, sendo cinco cursos com conceito 3 ( modal ) e um curso com conceito 4.
Outros dois cursos ficaram Sem Conceito ( SC ) e nenhum obteve conceito 1, 2 ou 5 nessa região.
Na região Centro Oeste, 25 dos 28 cursos participantes eram de IES privadas ( 89,3 \% em termos regionais ).
Desses, 13 cursos receberam conceito 3, o valor modal.
Dos demais cursos, dois receberam o conceito 2 ; oito receberam o conceito 4 ; e um curso recebeu o conceito 5.
Um curso ficou Sem Conceito ( SC ) e nenhum recebeu conceito 1.
Quanto aos 3 cursos oferecidos por instituições públicas na região Centro Oeste ( 10,7 \% do total regional ), foram atribuídos conceitos 3, 4 e 5 para cada um deles.
Nenhum curso ficou Sem Conceito ( SC ) nessa região.
Também na região Centro Oeste, a maioria dos cursos foi presencial : 27 ( 96,4 \% ).
Para esses, a moda foi o conceito 3, com 13 cursos ( 48,1 \% ).
Os demais cursos presenciais da região foram \textbf{alocados} nos conceitos 2 ( dois cursos, 7,4 \% ), 4 ( nove cursos, 33,3 \% ) e 5 ( dois cursos, 7,4 \% ).
Um curso presencial da região ficou Sem Conceito ( SC ) e nenhum obteve conceito 1.
Para o único curso a distância nessa região, o conceito atribuído foi 3.
À título de síntese, no Gráfico 6.2, apresenta -se a distribuição do Conceito Enade, segundo a categoria administrativa da IES.
Nesse gráfico, como dito anteriormente, poligonais mais à esquerda representam Conceito Enade mais baixos, e poligonais mais à direita, Conceito Enade mais altos.
É importante observar que, para cada poligonal presente nesse gráfico, a porcentagem constante no eixo y se refere ao universo de cursos oferecidos por instituições públicas ( poligonal vermelha ) ou privadas ( poligonal azul ), respectivamente. 6.3 CONCEITOS POR ORGANIZAÇÃO ACADÊMICA E GRANDE REGIÃO Na Tabela 6.3, encontra -se a distribuição dos conceitos atribuídos aos cursos participantes do Enade/2022, na área de Publicidade e Propaganda, por organização acadêmica, segundo as grandes regiões brasileiras.
Como mostram os dados da Tabela 6.3, dos 345 cursos de Publicidade e Propaganda participantes, 150 eram oferecidos em universidades, 121 em centros universitários e 74 em faculdades. nenhum curso foi oferecido por CEFET/IF.
Esta distribuição corresponde a, respectivamente, 43,5 \%, 35,1 \% e 21,4 \% dos cursos.
De acordo com os dados apresentados na mesma tabela, de todos os 29 cursos avaliados com o conceito 5, 14 eram vinculados a universidades.
Os cursos vinculados a esse tipo de organização acadêmica tiveram o conceito 3 como conceito modal, atribuído a 65 cursos ( 43,3 \% ).
Os demais cursos ofertados por universidades receberam os conceitos 1 ( três cursos ), 2 ( 23 cursos ), 4 ( 42 cursos ) e conceito 5 ( 14 cursos, como já mencionado ).
Três cursos ficaram Sem Conceito ( SC ).
Dos cursos vinculados a centros universitários, o conceito modal também foi o 3, atribuído a 61 cursos ( 50,4 \% ).
Onze cursos vinculados a esse tipo de organização acadêmica ficaram Sem Conceito ( SC ), e os outros cursos receberam os conceitos 1 ( dois cursos ), 2 ( 13 cursos ), 4 ( 25 cursos ) e 5 ( nove cursos ).
Dos 74 cursos mantidos por faculdades, 28 ( 37,8 \% ) receberam conceito 3, o conceito modal.
O segundo conceito mais frequentemente atribuído aos cursos oferecidos por essa categoria de IES foi o conceito 4, recebido por 18 cursos.
Na sequência, 11 cursos obtiveram conceito 2 ; seis cursos, o conceito 5 ; e dois cursos, o conceito 1.
Dos cursos oferecidos em Faculdades, nove ficaram Sem Conceito ( SC ).
Considerando -se, separadamente, as regiões brasileiras, verifica -se que, na região Norte, as universidades participaram com 5 cursos avaliados, aos quais foram atribuídos os seguintes conceitos : 2, a um curso ; 3, o conceito modal, a três cursos e conceito 4, também a um curso.
Nenhum curso ficou Sem Conceito ( SC ) ou com conceito 1 e 5 nessa combinação de categoria e região.
Os centros universitários da região Norte foram representados por seis cursos, dos quais três obtiveram conceito 3, o valor modal, e um curso, o conceito 5.
Os conceitos 1, 2 e 4 não foram atribuídos a nenhum curso.
Dois cursos ficaram Sem Conceito ( SC ).
As Faculdades participaram com 3 cursos na região Norte, os quais receberam os conceitos 3 ( dois cursos ) e 4 ( um curso ).
Nenhum curso ficou Sem Conceito ( SC ), e nenhum curso oferecido em faculdades da região obteve o conceito 1, 2 ou 5.
Na região Nordeste, as universidades participaram com 16 dos 46 cursos da área de Publicidade e Propaganda oferecidos.
O conceito modal foi 4, atribuído a seis cursos.
Outros quatro cursos obtiveram o conceito 3 ; três cursos, o conceito 2 ; e outros dois, o conceito 5.
Nenhum curso oferecido por universidades no Nordeste recebeu o conceito 1 e apenas um ficou Sem Conceito ( SC ).
Os Centros Universitários contaram com 18 cursos participantes na região Nordeste, 10 dos quais receberam o conceito modal 3.
Para os demais cursos, foram atribuídos os conceitos 2 ( três cursos ), 4 ( um curso ) e 5 ( um curso ).
Três cursos oferecidos por centros universitários no Nordeste ficaram Sem Conceito ( SC ) ; nenhum recebeu o conceito 1.
As faculdades foram representadas por 12 cursos na região Nordeste, quatro dos quais receberam conceito 2 e quatro, o conceito 3, os conceitos modais.
Aos demais cursos foram atribuídos os conceitos 1 ( um curso ), 4 ( dois cursos ) e 5 ( um curso ).
Nenhum curso ficou Sem Conceito ( SC ).
Na região Sudeste, as universidades concentraram 86 dos 188 cursos de Publicidade e Propaganda da região.
Dos cursos oferecidos em IES com este tipo de organização acadêmica, nessa região, o conceito modal foi 3, atribuído a 42 cursos.
Os demais cursos receberam os conceitos 1 ( três cursos ), 2 ( 18 cursos ), 4 ( 16 cursos ) e 5 ( seis cursos ).
Um curso ficou Sem Conceito ( SC ).
Os Centros Universitários participaram com 58 cursos na região Sudeste, dos quais 30 obtiveram o conceito modal 3.
O conceito 1 foi atribuído a dois cursos ; o conceito 2, a cinco cursos ; o conceito 4, a 15 cursos ; e o conceito 5, a três cursos.
Três cursos ficaram Sem Conceito ( SC ) nessa combinação de organização acadêmica e região.
Já as faculdades foram representadas por 44 cursos na região Sudeste, dos quais sete ficaram Sem Conceito ( SC ).
O conceito modal 3 foi atribuído a 14 cursos.
Aos demais cursos foram atribuídos os conceitos 1 ( um curso ), 2 ( cinco cursos ), 4 ( 13 cursos ) e 5 ( quatro cursos ).
Dos 69 cursos da região Sul, 34 eram mantidos por Universidades, para os quais o conceito modal foi 4, atribuído a 14 cursos.
Os demais cursos receberam os conceitos 2 ( um curso ), 3 ( 13 cursos ), 5 ( cinco cursos ) ; nenhum curso recebeu o conceito 1.
Em IES com esse tipo de organização acadêmica, um curso da região Sul ficou Sem Conceito ( SC ).
Dos 26 cursos participantes de centros universitários da região Sul, dois ficaram Sem Conceito ( SC ), e 10 cursos receberam conceito 3, o valor modal.
Os conceitos restantes ficaram assim distribuídos : cinco cursos ( conceito 2 ) ; seis cursos ( conceito 4 ) ; e três cursos ( conceito 5 ).
Quanto aos 9 cursos vinculados a faculdades na região Sul, cinco cursos receberam o conceito 3, valor modal, e dois cursos ficaram Sem Conceito ( SC ).
Os dois cursos restantes receberam os conceitos 4 e 5 ( um curso cada ).
Na região Centro-Oeste, 9 dos 28 cursos oferecidos eram mantidos por universidades.
Para esse tipo de organização acadêmica, o conceito modal foi 4, atribuído a cinco cursos ; e, nenhum curso da região ficou Sem Conceito ( SC ) ou obteve conceitos 1 ou 2.
Os demais cursos receberam os conceitos 3 ( três cursos ) e 5 ( um curso ).
Os centros universitários da região Centro-Oeste contaram com 13 cursos : oito cursos foram avaliados com o conceito 3 ( modal ) ; três, o conceito 4 e apenas um curso obteve o conceito 5.
Um outro curso ficou Sem Conceito ( SC ).
Dos 6 cursos oferecidos em Faculdades na região Centro-Oeste, três cursos receberam o conceito 3 ( modal ) ; dois cursos, o conceito 2 ; e um curso, o conceito 4.
Nessa região, nenhum curso oferecido em faculdades ficou Sem Conceito ( SC ) ou obteve os conceitos 1 e 5.
De maneira semelhante à forma como se procedeu, em relação à seção anterior, encerra -se com o Gráfico 6.3, que apresenta a distribuição do Conceito Enade dos cursos de Publicidade e Propaganda segundo a organização acadêmica da IES.
No Gráfico 6.3, apresenta -se a distribuição do Conceito Enade segundo a organização acadêmica da IES.
A leitura desse último gráfico é semelhante aos anteriores, seguindo -se a lógica de que poligonais deslocadas para a esquerda correspondem a Conceito Enade mais baixos ( 1, 2 e 3 ) e poligonais deslocadas para a direita correspondem a Conceito Enade mais altos.

% matched lemmas: alocar
\end{textsample}
